\chapter{Idiomatic Code Through the Rules}
\label{sec:idioms}

This chapter answers the question the thesis is organised around: \emph{how should the language and
the checker be designed so that idiomatic code almost never hits the ownership error?} It enumerates
the idioms that must always be accepted, shows for each the analysis precision that accepts it and the
language or library choice that makes it provable, lists what is still rejected, and re-assesses the
two applications we examined.

\section{The question, and what counts as idiomatic}
\label{sec:idioms-question}

An \emph{idiom} here is a way of writing a list operation that beni's language and library make the most
natural one and that an Elm programmer would write unprompted. Elm's own idioms carry over almost
verbatim: \code{x :: acc} is written \code{[ x, …acc ]}, and the formatter's migration does the rewrite.
Because list construction consumes its base (\S\ref{sec:model-demands}), nearly every idiom that builds
a list from another is a demand, and the checker meets demands everywhere. A text search over beni's
core library, its platforms, the compiler's test programs and the example applications---approximate,
with the method in Appendix~\ref{app:method}, and not the output of the analysis---finds about 112 list
expressions that prepend with a spread, 30 that append with one, 88 of the form \code{e ++ [ ⋯ ]}, 16
fold callbacks of the form \code{λx acc → [ x, …acc ]}, and 56 \code{update} arms, in 37 files, that
extend a model list with a spread or \code{++}. The core library's own \code{List.push} calls, eight of
them, are all accumulators.

The design target follows. For the idioms below, the checker should produce \emph{no} error at all;
what it rejects should be a short, named list of forms, each either real sharing that the program would
otherwise get wrong or a stated limit, each with a one-line fix in its message; and the frequency of the
error on idiomatic code is the main number the evaluation measures (\S\ref{sec:eval-decisions}).

\section{The idioms that must be accepted}
\label{sec:idioms-list}

\paragraph{I1. An accumulator in a loop.} From the core library and a test program:
\begin{Verbatim}
combineHelp rest (List.push values value)              -- Result.combine's helper
prependAll (n - 1) [ size - n + 1, …acc ]               -- prepend, then reverse
growPlus (n - 1) (acc ++ [ size - n + 1 ])
\end{Verbatim}
\emph{What accepts it.} The accumulator is a parameter; the construction consumes it; the tail call
rebinds it, so its old value is dead after the jump (\S\ref{sec:model-liveness}). The function's summary
says $\mathit{use}(\mathit{acc}) = \{\mathsf{R}, \mathsf{A}, \mathsf{W}\}$ and $\mathit{res} = \{\pi_2\}$
after one round of the group fixpoint, and the caller passes \code{[]}, a fresh array
(\S\ref{sec:hard-recursion}). \emph{What makes it provable.} Construction consumes its base, so the
idiom needs no writer's name; strict evaluation and the tail-call rebinding fix when the old value dies.

\paragraph{I2. A fold.} From the core library:
\begin{Verbatim}
keys dict = foldl dict [] λkey _ keyList → List.push keyList key
foldr d [] (λkey _ names → [ key, …names ])             -- the documented example of Dict.foldr
values fibers = List.foldr fibers (Just []) λf acc →
    case acc of
        Just rest → ⋯ Just [ v, …rest ] ⋯
\end{Verbatim}
\emph{What accepts it.} The callback writes its own accumulator parameter, so the edge of
\S\ref{sec:inf-classes} requires the fold to supply that parameter owned, and the fold's row says it
does: the accumulator is the fold's own parameter on the first call and the callback's previous result
afterwards, neither held by anything else. In the third example the accumulator is wrapped in
\code{Just}; field-sensitive paths through constructors give the callback $\mathit{use}(\pi_2.\mathit{Just}\#0) \ni \mathsf{W}$,
and the supply is owned at every path. \emph{What makes it provable.} The folds' rows (asserted or
inferred) supply the accumulator owned; saturated calls put every callback in front of the checker as a
lambda; there is no point-free plumbing through which its class could be lost.

\paragraph{I3. A recursive builder.} From the compiler's test programs:
\begin{Verbatim}
mapRec xs f = case xs of
    [] → []
    [ x, …rest ] → [ f x, …mapRec rest f ]
\end{Verbatim}
\emph{What accepts it.} The base is the recursive call's result, whose provenance is $\fresh$ (the base
case's \code{[]}) after the group fixpoint; the pattern only reads, and \code{xs} is dead after the match.
A builder whose base case returns a parameter, such as an \code{appendTo} that returns its second
argument, consumes that parameter, and its callers must not use it afterwards (\S\ref{sec:model-demands}).
\emph{What makes it provable.} Patterns only read and bind views; the backend's building loop keeps the
source's meaning.

\paragraph{I4. A list built by \code{map} or \code{filter}, then extended; a pipeline.}
\begin{Verbatim}
visible = model.todos ▷ List.filter isOpen ▷ List.map label
[ …List.filter xs p, y ]
\end{Verbatim}
\emph{What accepts it.} \code{map} and \code{filter} are not demands; their results are $\{\fresh,
\pi_1\}$, the $\pi_1$ because of the identity promises the language keeps (\S\ref{sec:model-demands}). A
construction on such a result consumes the source too, so the source must be dead there; in an
\code{update} arm and in a local pipeline it is. \emph{What makes it provable.} The pipe is rewritten to
a call before anything runs, so there is no plumbing function in which to lose precision
(\S\ref{sec:hard-pipelines}); the identity promises cost precision only where the source is used again.

\paragraph{I5. An \code{update} that extends or edits a model list.} From the compiler's test programs
and the two applications:
\begin{Verbatim}
{ model | seen = [ …model.seen, k ] }
( { model | errors = model.errors ++ Api.errorMessages error }, Cmd.none )
changed { model | draft = "", nextId = model.nextId + 1 } (model.todos ++ [ new ])
Undo → case model.history of
    [ previous, …rest ] →
        { model | notes = previous, history = rest, future = [ model.notes, …model.future ] }
\end{Verbatim}
\emph{What accepts it.} The entry protocol (\S\ref{sec:hard-tea}) hands \code{update} the model owned
(\Obl{P1}) and computes, from the summaries of \code{init} and \code{update}, which cells the model's
paths may hold at entry; a path is editable when they are fresh, unaliased with other paths, and neither
escaped nor shared. Inside the arm, the record update replaces \code{seen} and shares the other fields, so
per-path liveness finds no later use of the old \code{model.seen}. In the third example the first
argument, a record still holding the old \code{todos}, is pending while the second argument consumes it;
it is not a use, because \code{changed} never reads its first argument's \code{.todos}
($\mathit{use} = \unused$), so the pending holder is dead. In the undo arm, \code{history} becomes a view
of the old history and \code{future} is consumed; the old model is dead. \emph{What makes it provable.}
Four choices, three of them outside the checker: \code{init} is a function called per mount, and a
top-level list is shared, so the model starts with cells nothing else holds; the runtime hands the model
over; a renderer does not skip work by the identity of a list the program edits in place (\Obl{P7}); and
the program persists and reports through commands that capture values rather than lists (below).

\paragraph{I6. An edit inside a container.}
\begin{Verbatim}
Dict.update d (key x) (λm → Just [ x, …Maybe.withDefault m [] ])     -- grouping
List.update rows i (λr → { r | tags = [ …r.tags, t ] })
\end{Verbatim}
\emph{What accepts it.} The container's contents are exclusive (\S\ref{sec:model-excl}), because they were
built by stores of fresh, distinct values, and the container function supplies the element owned. For a
dictionary the owned supply is asserted, and requires its keys and values to be disjoint as well
(\S\ref{sec:hard-containers}). \emph{What makes it provable.} The library offers the in-place
\code{update} at a key or an index, which is also Elm's idiom for grouping; a lookup followed by an
insert is real sharing (\S\ref{sec:idioms-rejected}).

\paragraph{I7. Several accumulators in a tuple or record.} From the core library and an idiomatic fold:
\begin{Verbatim}
partitionHelp rest errors (List.push values value)                     -- Result.partition
List.foldl xs ( [], [] ) λx ( yes, no ) →
    if p x then ( [ …yes, x ], no ) else ( yes, [ …no, x ] )
\end{Verbatim}
\emph{What accepts it.} Paths through tuples and records keep the accumulators apart; the callback's
summary writes $\pi_2.1$ and $\pi_2.2$ separately; the fold supplies each owned. \emph{What makes it
provable.} Field-sensitive paths and per-path supply; nothing in the idiom names a writer.

\section{The precision each idiom needs}
\label{sec:idioms-precision}

Table~\ref{tab:idioms} sums up which part of the analysis each idiom rests on. Every entry is a part of
the design of Chapters~\ref{sec:model} and~\ref{sec:inference}; none is an extension still to be
designed. The table also shows what an implementation must get right first: per-path liveness, callback
classes with owned supply, and the entry protocol carry almost every idiom.

\begin{table}[t]
\centering\small
\begin{tabularx}{\linewidth}{@{}>{\raggedright\arraybackslash}p{0.42\linewidth}>{\raggedright\arraybackslash}p{0.2\linewidth}X@{}}
\toprule
Part of the analysis & Idioms & Where \\
\midrule
per-path liveness over the evaluation-order normal form & I1, I5, I7 & \S\ref{sec:model-liveness} \\
the rebinding of a self tail call & I1 & \S\ref{sec:model-liveness} \\
the group fixpoint from optimistic summaries & I1, I3 & \S\ref{sec:inf-recursion} \\
callback classes: owned supply, provenance & I2, I6, I7 & \S\ref{sec:inf-classes} \\
paths through constructors, tuples and records & I2, I5, I7 & \S\ref{sec:model-cells} \\
$\unused$ parameter paths; pending operands & I5 & \S\ref{sec:model-summary}, \S\ref{sec:model-unique} \\
results naming both outcomes of an identity promise & I4 & \S\ref{sec:model-demands} \\
the entry protocol and its fixpoint over $\fresh$, $\esc$, $\shared$ & I5 & \S\ref{sec:hard-tea} \\
exclusive contents; a dictionary's split step & I6 & \S\ref{sec:model-excl} \\
\bottomrule
\end{tabularx}
\caption{The parts of the analysis that accept each idiom.}
\label{tab:idioms}
\end{table}

The language and library choices that make the idioms provable are, in order of reach:
\begin{enumerate}
\item \emph{Construction consumes its base, by position.} Every idiom above runs in place without a new
name, and the programmer chooses which operand is reused by where it is written.
\item \emph{Strict left-to-right evaluation and saturated calls.} They make ``before'' a fact of the text
and every closure a lambda in front of the checker.
\item \emph{Patterns only read.} Walking a list never consumes it.
\item \emph{The folds and container updates supply their callback owned}, by their rows.
\item \emph{\code{init} is a function, and top-level lists are shared.} A model starts with cells no one
else holds, and a constant is never edited behind its readers.
\item \emph{The runtime hands the model over, and a renderer does not skip work by the identity of a list
the program edits} (\Obl{P1}, \Obl{P7}).
\item \emph{Commands and messages carry values.} A command that persists or reports a list should
capture the value it needs---the encoded text, the count---rather than a closure over the list, as Elm's
own \code{Http.jsonBody} encodes when the command is made; and a view event should carry an identifier
rather than a model list. Both are library and teaching choices, not analysis.
\item \emph{\code{eq} and \code{compare} only read}, so sorting and dictionaries never write behind the
program's back.
\end{enumerate}

\section{What is still rejected}
\label{sec:idioms-rejected}

Table~\ref{tab:rejected} lists every form we know that the rules reject in code written the way the
language invites, with the reason, the fix the message gives, and the choice that would make the form
rarer. Rows marked \emph{real} are real sharing: under in-place semantics the program would print
something different. Rows marked \emph{limit} are the checker's or the platform's blindness. The
evaluation counts each row separately.

{\small
\begin{longtable}{@{}>{\raggedright\arraybackslash}p{0.24\linewidth}>{\raggedright\arraybackslash}p{0.29\linewidth}>{\raggedright\arraybackslash}p{0.12\linewidth}>{\raggedright\arraybackslash}p{0.26\linewidth}@{}}
\caption{Forms the rules reject in idiomatic-looking code.}
\label{tab:rejected}\\
\toprule
Form & Example & Kind & Fix; what makes it rarer \\
\midrule
\endfirsthead
\toprule
Form & Example & Kind & Fix; what makes it rarer \\
\midrule
\endhead
a version kept and then extended & a persistent stack whose old versions are read; \code{a = base ++ [ 1 ]; b = base ++ [ 2 ]} & real & \code{List.copy} at the kept version, $O(n)$ and written; a persistent-vector library for programs that keep many versions \\
the base read again in the same expression & \code{\{ model | items = [ …model.items, x ], count = List.length model.items \}}; \code{[ …xs, …xs ]} & real & compute the read first; \code{List.copy} \\
a list captured by a command, then extended & TodoMVC's \code{save todos = Cmd.do λ⊤ → ⋯ (encode todos)} & real, possibly & compute before the command; command constructors that take values \\
a top-level list as a base & \code{[ x, …defaults ]}; \code{init ⊤ = \{ todos = seed \}} & real & build in a function; \code{List.copy}; the reading form \code{[ x ] ++ defaults} \\
a model list as a base in \code{view}, a derived value or \code{subscriptions} & \code{[ …model.pinned, …model.recent ]} in \code{view} & real & \code{List.concat [ model.pinned, model.recent ]}, which builds \\
a once closure passed to a $\mathit{many}$ position & \code{List.map xs (λx → [ …acc, x ])} & real & a fold \\
a lookup, an edit and a re-insert & \code{Dict.get}, then \code{[ …group, x ]}, then \code{Dict.insert} & real & \code{Dict.update} \\
a nested edit in a container not known exclusive & elements from \code{List.repeat}, an undeclared decoder, JavaScript & limit or real & \code{List.copy} of the inner list; decoders declare exclusive output \\
a list from a \code{Ref} or from JavaScript & \code{Ref.update r (λxs → [ …xs, x ])} & limit & copy where it arrives \\
a list a view event carries & \code{onClick=\{Keep model.rows\}} anywhere in \code{view} & limit & send an identifier \\
a rendered list on a renderer that skips by identity & any of I5 on the template runtime without \Obl{P7} & limit (platform) & the platform adopts \Obl{P7} \\
\bottomrule
\end{longtable}
}

Three rows deserve comment. The first reverses part of a decision the language specification records:
that no way of building a list by prepending is quadratic, a persistent stack included
(\S\ref{sec:bg-lists}). Every accumulator, fold, recursive builder and TEA list keeps linear prepending
under the rules, in place. A persistent stack whose old versions are read later does not: it is rejected
until each kept version is copied, which costs $O(n)$ per kept version, written in the source. The test
program that guards the decision builds four accumulators, which are accepted in place, and churns a front
stack that keeps every 10\,000th version, which is rejected at its prepend: a kept version may share the array, and a prepend after a pop
would overwrite a slot that the kept version still shows. With \code{List.copy} at the version it keeps it is accepted, with ten copies
of at most 100\,000 elements. The trade is the one the rule exists to make: a version costs what it
costs, visibly.

The third row, commands, is the largest source of error in the two applications, and the cheapest to
remove: the fix is one binding before the command. The last row is not about programs at all; it is the
renderer, and it decides whether TEA idioms are accepted on beni's current browser platform
(\S\ref{sec:hard-tea}).

\section{The two applications, re-assessed}
\label{sec:disc-apps}

We examined two applications written in beni: TodoMVC, in five programs---the full one among the
compiler's test programs, and four variants used to measure output size---and Conduit, the RealWorld
example application. Both run on beni's TEA platform over the template runtime. What follows is our
reading of their sources, traced through the rules by hand; it is not the output of an analysis, which
does not exist yet.

\paragraph{TodoMVC.} Each program has one construction that is a demand: the Enter arm's
\code{changed \{ model | draft = "", ⋯ \} (model.todos ++ [ new ])}, which appends in place. Its other arms
build with \code{List.map} and \code{List.filter}, which are not demands. Four facts decide the outcome.
\begin{itemize}
\item \code{init} is a value today; under the rules it becomes \code{init ⊤ = ⋯}, a mechanical change. It
decodes the stored list or starts from \code{[]}, so the model's \code{todos} is fresh at the first entry.
\item In the Enter arm the pending first argument still holds the old \code{todos}, but \code{changed}'s
summary is $\unused$ at $\pi_1.\mathit{todos}$, so the append is not seen through it.
\item The full program and two of the variants persist with
\code{save todos = Cmd.do λ⊤ → saved (Storage.set Local key (encode todos))}, called by every arm that
changes the list. The command's body captures the arm's new list, so \code{update}'s summary says the
model's \code{todos} may be $\esc_k$, and at every later entry the path is escaped: the append is rejected
as possible sharing, on every platform. Under value semantics the command encodes the list as the arm left it; in place, a
command that ran after a later append would encode the appended list, so the rejection is right. The
fix is one binding before the command:
\begin{Verbatim}
save todos =
    text = encode todos
    Cmd.do λ⊤ → saved (Storage.set Local key text)
\end{Verbatim}
after which the command captures a string.
\item The page renders \code{<For each=\{List.filter model.todos (visible model.filter \_)\}>}. By the
identity promise, that list may be \code{model.todos} itself, and the template runtime skips the walk when
it receives the same list as last time. An in-place append keeps the list's identity, so the new row would
not be drawn. On the template runtime as it is, the append is therefore rejected; with \Obl{P7} it is
accepted, and that \code{<For>} walks its rows on every render.
\end{itemize}
So, as written, no TodoMVC program is accepted. With \code{init} a function and, where it persists, the
one-line change to \code{save}, all five are accepted on the direct platform and, if the template runtime
adopts \Obl{P7}, on their own platform. What they gain is an append that writes one slot instead of
copying the list, which at the sizes TodoMVC handles is invisible, and the smaller list runtime.

\paragraph{Conduit.} Its \code{init} is already a function of the URL and the navigation key. It has
thirteen \code{++} on lists. Six have a left operand built in the same function---a literal such as
\code{[ "api" ] ++ paths}, or a local list in a view helper---and consume only that. Seven are
\code{model.errors ++ Api.errorMessages error} in the arms of three page modules: in-place appends onto
the page's error list, which each page's \code{init} creates as \code{[]}, which no command captures, by
our reading, and which the page renders with a positional \code{<For>}. The derivation needs paths through
the constructor that holds the current page in the top-level model and through the placeholder lambdas
that wrap a page's model back into it; both are ordinary paths and classes. On the direct platform, by
our reading, all seven appends are accepted. On the template runtime they are rejected unless it adopts
\Obl{P7}, for the same reason as TodoMVC's.

\paragraph{What the applications say.} Under construction-as-consumption the applications do edit their
lists in place, in the idiom I5, and the rule's outcome for them turns on two things that are not in the
checker: a one-line persistence change in one of them, and the template runtime's identity skip in both.
The second is a platform decision with a cost per render (\S\ref{sec:hard-tea}); the evaluation measures
that cost against the reference renderer before it is adopted. The applications are small and their
lists short; they test whether the rule is livable, not whether it is fast.
